\documentclass[aspectratio=169]{beamer}

% Tema UFS — Universidade Federal de Sergipe
\usetheme{UFS}

% Pacotes base
\usepackage[utf8]{inputenc}
\usepackage[portuguese]{babel}
\usepackage{amsmath,amssymb}
\usepackage{graphicx}
\usepackage{booktabs}
\usepackage{xcolor}

% TikZ e bibliotecas
\usepackage{tikz}
\usetikzlibrary{shapes.geometric, arrows.meta, positioning, matrix, fit, calc,
  backgrounds, decorations.pathreplacing, shadows, patterns}

% Listagens — paleta VS Code Light+
\usepackage{listings}
\definecolor{bgcolor}{HTML}{FFFFFF}
\definecolor{linecolor}{HTML}{DDDDDD}
\definecolor{numcolor}{HTML}{999999}
\definecolor{kwcolor}{HTML}{0000FF}
\definecolor{strcolor}{HTML}{A31515}
\definecolor{cmtcolor}{HTML}{008000}

\lstdefinestyle{ufsStyle}{
  backgroundcolor=\color{bgcolor},
  basicstyle=\ttfamily\footnotesize\color{black},
  keywordstyle=\color{kwcolor}\bfseries,
  stringstyle=\color{strcolor},
  commentstyle=\color{cmtcolor}\itshape,
  numberstyle=\tiny\color{numcolor},
  numbers=left, numbersep=12pt,
  xleftmargin=20pt, framexleftmargin=20pt,
  breaklines=true, tabsize=4,
  keepspaces=true, showspaces=false,
  showstringspaces=false, showtabs=false,
  frame=single, rulecolor=\color{linecolor},
  framesep=4pt, captionpos=b, upquote=true,
  columns=fullflexible,
}
\lstset{style=ufsStyle, language=C}

% --- Metadados ---
\title{Algoritmos e Estruturas de Dados II}
\subtitle{Corretude de algoritmos iterativos e invariantes de laço --- Parte 1}
\author{Prof.\ André Yoshiaki Kashiwabara}
\institute{DCOMP -- Universidade Federal de Sergipe\\
\texttt{andreyoshiaki@dcomp.ufs.br}}
\date{}

\begin{document}

\begin{frame}[plain]
  \titlepage
\end{frame}

\begin{frame}{O que você vai aprender hoje}
  \begin{center}
  \begin{tikzpicture}[
    item/.style={rectangle, draw=UFSBlue, fill=UFSBlue!10, rounded corners,
      minimum width=9.5cm, minimum height=0.75cm, font=\small}
  ]
    \node[item] (t1) at (0,0)    {Por que testar um algoritmo não basta para garantir que ele é correto};
    \node[item] (t2) at (0,-1.1) {O que é um \textbf{invariante de laço} e as três propriedades que o tornam uma prova};
    \node[item] (t3) at (0,-2.2) {Prova de corretude da \textbf{busca do máximo} em um vetor};
    \node[item] (t4) at (0,-3.3) {Prova de corretude da \textbf{ordenação por inserção}};
    \draw[->, thick, UFSBlue] (t1)--(t2);
    \draw[->, thick, UFSBlue] (t2)--(t3);
    \draw[->, thick, UFSBlue] (t3)--(t4);
  \end{tikzpicture}
  \end{center}
\end{frame}

\begin{frame}{Agenda}
  \tableofcontents
\end{frame}

% =====================================================================
\section{Corretude e invariantes de laço}
% =====================================================================

\begin{frame}{Por que precisamos provar corretude?}
  \begin{columns}[T]
    \begin{column}{0.5\textwidth}
      \textbf{Problema:}
      \begin{itemize}
        \item Testes mostram que o algoritmo funciona \emph{nos casos testados} --- nunca em todos.
        \item Um vetor de 20 inteiros de 32 bits tem mais entradas possíveis do que átomos observáveis conseguiríamos testar.
        \item Laços são a principal fonte de erros sutis: uma iteração a mais, uma a menos, um índice trocado.
      \end{itemize}
    \end{column}
    \begin{column}{0.5\textwidth}
      \textbf{Solução:}
      \begin{itemize}
        \item Um \emph{argumento matemático} que vale para \textbf{toda} entrada válida.
        \item Para algoritmos iterativos, a ferramenta central é o \textbf{invariante de laço}.
        \item A mesma ideia que usamos para projetar o laço serve para \emph{provar} que ele funciona.
      \end{itemize}
    \end{column}
  \end{columns}
\end{frame}

\begin{frame}{O que significa ``o algoritmo é correto''?}
  \begin{block}{Definição (corretude)}
    Um algoritmo é \textbf{correto} quando, para \emph{toda} entrada que satisfaz a
    \textbf{precondição}, ele \textbf{termina} e sua saída satisfaz a \textbf{poscondição}.
  \end{block}
  \begin{exampleblock}{Exemplo: busca do máximo}
    \textbf{Precondição:} vetor $A[0..n-1]$ com $n \geq 1$ elementos.\\
    \textbf{Poscondição:} o valor devolvido é $\max\{A[0], A[1], \ldots, A[n-1]\}$.
  \end{exampleblock}
  \begin{itemize}
    \item Corretude \emph{parcial}: \textbf{se} terminar, a resposta é certa.
    \item Corretude \emph{total}: parcial $+$ garantia de \textbf{término}.
  \end{itemize}
\end{frame}

\begin{frame}{Invariante de laço: definição}
  \begin{block}{Definição}
    Um \textbf{invariante de laço} é uma afirmação sobre as variáveis do programa que é
    \textbf{verdadeira no início de cada iteração} do laço.
  \end{block}
  \begin{exampleblock}{Intuição}
    É a ``promessa'' que o laço mantém enquanto trabalha: descreve \emph{o que já foi
    conquistado} até aquele ponto. Ao final, a promessa acumulada deve implicar o
    resultado desejado.
  \end{exampleblock}
\end{frame}

\begin{frame}{As três propriedades de uma prova por invariante}
  \begin{center}
  \begin{tikzpicture}[
    fase/.style={rectangle, draw=UFSBlue, fill=UFSBlue!10, rounded corners,
      minimum width=3.6cm, minimum height=1.2cm, font=\small, align=center},
    seta/.style={->, thick, UFSBlue}
  ]
    \node[fase] (ini) at (0,0)
      {\textbf{Inicialização}\\ vale \emph{antes} da\\ primeira iteração};
    \node[fase] (man) at (5,0)
      {\textbf{Manutenção}\\ cada iteração\\ preserva o invariante};
    \node[fase] (ter) at (10,0)
      {\textbf{Término}\\ ao sair, invariante\\ $\Rightarrow$ poscondição};
    \draw[seta] (ini)--(man);
    \draw[seta] (man)--(ter);
    \draw[seta] (man) to[out=140, in=40, looseness=4] (man);
  \end{tikzpicture}
  \end{center}
  \begin{itemize}
    \item É o esquema da \textbf{indução matemática}: inicialização é o caso base,
      manutenção é o passo indutivo --- e a ``indução'' para quando o laço termina.
    \item Para corretude \emph{total}, mostramos também que o laço \textbf{termina}
      (ex.: uma quantidade inteira não negativa que decresce a cada iteração).
  \end{itemize}
\end{frame}

\begin{frame}{Recapitulando: corretude e invariantes}
  \begin{block}{O que aprendemos}
    \begin{itemize}
      \item Testar mostra presença de acertos, não ausência de erros: corretude exige prova.
      \item Corretude = precondição $\rightarrow$ poscondição $+$ término.
      \item Invariante de laço: afirmação verdadeira no início de toda iteração.
      \item Prova em três passos: inicialização, manutenção e término.
    \end{itemize}
  \end{block}
\end{frame}

% =====================================================================
\section{Primeiro algoritmo: busca do máximo}
% =====================================================================

\begin{frame}{Por que começar pela busca do máximo?}
  \begin{columns}[T]
    \begin{column}{0.5\textwidth}
      \textbf{Problema:}\\[2pt]
      Dado um vetor $A[0..n-1]$ com $n \geq 1$, encontrar o valor máximo.
      \begin{itemize}
        \item Algoritmo de uma linha de raciocínio --- ideal para ver a \emph{mecânica} da prova sem distrações.
      \end{itemize}
    \end{column}
    \begin{column}{0.5\textwidth}
      \textbf{Solução:}\\[2pt]
      Percorrer o vetor guardando o maior valor \emph{já visto}.
      \begin{itemize}
        \item ``O maior já visto'' é exatamente o invariante --- a frase informal \emph{é} a afirmação formal.
      \end{itemize}
    \end{column}
  \end{columns}
\end{frame}

\begin{frame}[fragile]{Busca do máximo: o algoritmo}
\begin{lstlisting}
int maximo(int A[], int n) {   /* pre: n >= 1            */
    int max = A[0];
    for (int i = 1; i < n; i++) {
        /* invariante: max == maior valor de A[0..i-1]   */
        if (A[i] > max)
            max = A[i];
    }
    return max;                /* pos: max(A[0..n-1])    */
}
\end{lstlisting}
  \begin{block}{Invariante}
    No início da iteração de índice $i$:\quad $\texttt{max} = \max\{A[0], \ldots, A[i-1]\}$.
  \end{block}
\end{frame}

\begin{frame}{Busca do máximo: o invariante em figura}
  \begin{center}
  \begin{tikzpicture}[
    cel/.style={draw, minimum width=0.95cm, minimum height=0.8cm, font=\small},
    visto/.style={cel, fill=UFSBlue!20},
    atual/.style={cel, fill=UFSYellow!50},
    resto/.style={cel, fill=UFSLightGray}
  ]
    \node[visto] (a0) at (0,0)    {3};
    \node[visto] (a1) at (0.95,0) {9};
    \node[visto] (a2) at (1.9,0)  {4};
    \node[atual] (a3) at (2.85,0) {7};
    \node[resto] (a4) at (3.8,0)  {1};
    \node[resto] (a5) at (4.75,0) {8};

    \foreach \i/\n in {0/a0,1/a1,2/a2,3/a3,4/a4,5/a5}
      \node[below=1pt of \n, font=\scriptsize, color=UFSGray] {\i};

    \draw[decorate, decoration={brace, amplitude=5pt}, thick, UFSBlue]
      (-0.47,0.5) -- (2.37,0.5)
      node[midway, above=6pt, font=\small, color=UFSBlue]
      {já visto: $\texttt{max} = 9$};
    \node[above=14pt of a3, font=\small] {$i = 3$};
    \draw[->, thick] ($(a3.north)+(0,0.35)$) -- (a3.north);
  \end{tikzpicture}
  \end{center}
  \begin{itemize}
    \item Azul: prefixo $A[0..i-1]$ já examinado --- o invariante fala \textbf{só} sobre ele.
    \item Amarelo: elemento da iteração atual; cinza: ainda não examinado.
    \item A iteração ``incorpora'' $A[i]$ ao prefixo e restaura o invariante para $i+1$.
  \end{itemize}
\end{frame}

\begin{frame}{Busca do máximo: prova de corretude}
  \begin{itemize}
    \item \textbf{Inicialização} ($i = 1$): $\texttt{max} = A[0] = \max\{A[0]\}$
      --- o invariante vale para o prefixo de tamanho 1. \checkmark
    \item \textbf{Manutenção}: suponha $\texttt{max} = \max\{A[0..i-1]\}$ no início da iteração $i$.
      \begin{itemize}
        \item Se $A[i] > \texttt{max}$: novo $\texttt{max} = A[i] = \max\{A[0..i]\}$.
        \item Senão: $A[i] \leq \texttt{max}$, logo $\texttt{max}$ já é $\max\{A[0..i]\}$.
      \end{itemize}
      Nos dois casos o invariante vale para $i+1$. \checkmark
    \item \textbf{Término}: o laço para com $i = n$; o invariante dá
      $\texttt{max} = \max\{A[0..n-1]\}$ --- exatamente a poscondição. \checkmark
    \item \textbf{O laço termina?} $n - i$ é inteiro, não negativo e decresce 1 por iteração. \checkmark
  \end{itemize}
\end{frame}

\begin{frame}{Exemplo: onde a prova pega um erro}
  \begin{exampleblock}{Versão com erro: e se o laço começasse com \texttt{max = 0}?}
    Com $\texttt{max} = 0$ e $i = 0$, a ``inicialização'' exigiria
    $0 = \max\{\,\} $ --- o prefixo vazio não tem máximo, e para o vetor
    $A = [-5, -2, -9]$ a função devolveria $0$, que nem está no vetor.
  \end{exampleblock}
  \begin{itemize}
    \item O passo de \textbf{inicialização} falha: a prova não sai do lugar --- e o erro
      aparece \emph{antes} de qualquer teste.
    \item Erro clássico de quem inicializa acumuladores ``no chute'': o invariante obriga
      a escolher $\texttt{max} = A[0]$ (e $i = 1$).
  \end{itemize}
\end{frame}

\begin{frame}{Recapitulando: busca do máximo}
  \begin{block}{O que aprendemos}
    \begin{itemize}
      \item O invariante formaliza a frase ``\texttt{max} é o maior já visto''.
      \item Os três passos da prova foram diretos: caso base no prefixo de tamanho 1,
        análise de dois casos na manutenção, poscondição no término.
      \item A prova diagnostica erros de inicialização que testes podem não revelar.
    \end{itemize}
  \end{block}
\end{frame}

% =====================================================================
\section{Segundo algoritmo: ordenação por inserção}
% =====================================================================

\begin{frame}{Por que a ordenação por inserção?}
  \begin{columns}[T]
    \begin{column}{0.5\textwidth}
      \textbf{Problema:}\\[2pt]
      Ordenar $A[0..n-1]$ em ordem não decrescente.
      \begin{itemize}
        \item Agora o invariante não fala de \emph{um valor}, mas de uma
          \emph{propriedade estrutural} de um trecho do vetor.
        \item Há dois laços aninhados --- um invariante para cada.
      \end{itemize}
    \end{column}
    \begin{column}{0.5\textwidth}
      \textbf{Solução:}\\[2pt]
      Como organizar cartas na mão: a cada passo, inserir a próxima carta na posição
      certa entre as que \emph{já estão ordenadas}.
      \begin{itemize}
        \item ``As cartas na mão estão ordenadas'' é o invariante do laço externo.
      \end{itemize}
    \end{column}
  \end{columns}
\end{frame}

\begin{frame}[fragile]{Ordenação por inserção: o algoritmo}
\begin{lstlisting}
void insercao(int A[], int n) {
    for (int i = 1; i < n; i++) {
        /* invariante: A[0..i-1] ordenado, com os        */
        /* mesmos elementos que estavam em A[0..i-1]     */
        int chave = A[i];
        int j = i - 1;
        while (j >= 0 && A[j] > chave) {
            A[j + 1] = A[j];
            j--;
        }
        A[j + 1] = chave;
    }
}
\end{lstlisting}
\end{frame}

\begin{frame}{O invariante do laço externo em figura}
  \begin{center}
  \begin{tikzpicture}[
    cel/.style={draw, minimum width=0.95cm, minimum height=0.8cm, font=\small},
    ord/.style={cel, fill=UFSBlue!20},
    atual/.style={cel, fill=UFSYellow!50},
    resto/.style={cel, fill=UFSLightGray}
  ]
    \node[ord]   (a0) at (0,0)    {2};
    \node[ord]   (a1) at (0.95,0) {5};
    \node[ord]   (a2) at (1.9,0)  {8};
    \node[atual] (a3) at (2.85,0) {4};
    \node[resto] (a4) at (3.8,0)  {1};
    \node[resto] (a5) at (4.75,0) {6};

    \foreach \i/\n in {0/a0,1/a1,2/a2,3/a3,4/a4,5/a5}
      \node[below=1pt of \n, font=\scriptsize, color=UFSGray] {\i};

    \draw[decorate, decoration={brace, amplitude=5pt}, thick, UFSBlue]
      (-0.47,0.5) -- (2.37,0.5)
      node[midway, above=6pt, font=\small, color=UFSBlue]
      {$A[0..i-1]$ ordenado};
    \node[above=14pt of a3, xshift=1.6cm, font=\small] (lbl) {$\texttt{chave} = A[i]$};
    \draw[->, thick] (lbl.south west) -- (a3.north);

    \draw[->, thick, UFSRed] (a3.south) to[out=-90, in=-90, looseness=1.6]
      node[below, font=\scriptsize, color=UFSRed] {inserir na posição certa}
      ($(a1.south)+(0.47,0)$);
  \end{tikzpicture}
  \end{center}
  \begin{block}{Invariante do laço externo}
    No início da iteração $i$: o trecho $A[0..i-1]$ está \textbf{ordenado} e contém
    \textbf{os mesmos elementos} que ocupavam $A[0..i-1]$ originalmente.
  \end{block}
\end{frame}

\begin{frame}{Por que ``os mesmos elementos''?}
  \begin{exampleblock}{Um algoritmo ``ordenado'' demais}
    \texttt{for (i = 0; i < n; i++) A[i] = i;} \quad deixa o vetor perfeitamente
    ordenado --- e completamente errado.
  \end{exampleblock}
  \begin{itemize}
    \item ``Estar ordenado'' sozinho é uma poscondição \textbf{fraca demais}: precisamos
      que o resultado seja uma \emph{permutação} da entrada.
    \item Por isso o invariante tem \textbf{duas} cláusulas: ordem $+$ mesmos elementos.
    \item Na inserção, só \emph{deslocamos} e \emph{recolocamos} a chave: nenhum valor é
      criado ou destruído --- a permutação é preservada.
  \end{itemize}
\end{frame}

\begin{frame}{Ordenação por inserção: prova (laço externo)}
  \begin{itemize}
    \item \textbf{Inicialização} ($i = 1$): $A[0..0]$ tem um único elemento ---
      ordenado e intacto. \checkmark
    \item \textbf{Manutenção}: o laço interno desloca uma posição à direita os elementos
      de $A[0..i-1]$ maiores que a chave e a insere logo antes deles.
      \begin{itemize}
        \item Invariante do laço interno: $A[j+2..i]$ contém os elementos deslocados,
          todos $> \texttt{chave}$, e $A[0..j]$ está intacto.
        \item Ao inserir a chave em $A[j+1]$: $A[0..i]$ fica ordenado, com os mesmos
          elementos de antes. \checkmark
      \end{itemize}
    \item \textbf{Término}: o laço para com $i = n$; o invariante dá $A[0..n-1]$
      ordenado e permutação da entrada --- a poscondição de ordenação. \checkmark
    \item \textbf{O laço termina?} Externo: $n - i$ decresce. Interno: $j + 1 \geq 0$
      decresce a cada iteração. \checkmark
  \end{itemize}
\end{frame}

\begin{frame}{Exemplo: rastreando o invariante em $A = [5, 2, 4, 1]$}
  \begin{center}
  \small
  \begin{tabular}{c l l}
    \toprule
    Início de $i$ & Vetor & Invariante: $A[0..i-1]$ ordenado? \\
    \midrule
    $i = 1$ & $[\mathbf{5} \mid 2, 4, 1]$ & $[5]$ \checkmark \\
    $i = 2$ & $[\mathbf{2, 5} \mid 4, 1]$ & $[2,5]$ \checkmark \\
    $i = 3$ & $[\mathbf{2, 4, 5} \mid 1]$ & $[2,4,5]$ \checkmark \\
    saída ($i = 4$) & $[\mathbf{1, 2, 4, 5}]$ & $[1,2,4,5]$ \checkmark\ $=$ poscondição \\
    \bottomrule
  \end{tabular}
  \end{center}
  \begin{itemize}
    \item A barra separa o prefixo ordenado (invariante) do restante do vetor.
    \item Rastrear o invariante em um exemplo \textbf{não é a prova} --- mas é o melhor
      jeito de \emph{descobrir} qual invariante provar.
  \end{itemize}
\end{frame}

\begin{frame}{Recapitulando: ordenação por inserção}
  \begin{block}{O que aprendemos}
    \begin{itemize}
      \item Invariantes podem afirmar propriedades estruturais (``prefixo ordenado''),
        não apenas valores.
      \item A poscondição de ordenação exige duas cláusulas: ordem e permutação da entrada.
      \item Laços aninhados: cada laço tem seu invariante, e o do interno sustenta a
        manutenção do externo.
    \end{itemize}
  \end{block}
\end{frame}

% =====================================================================
\section*{Encerramento}
% =====================================================================

\begin{frame}{Resumo da aula e o que vem a seguir}
  \begin{block}{Hoje}
    \begin{itemize}
      \item Invariante de laço $+$ três passos (inicialização, manutenção, término)
        $=$ prova de corretude de algoritmos iterativos.
      \item Busca do máximo: invariante sobre um valor acumulado.
      \item Ordenação por inserção: invariante estrutural com laços aninhados.
    \end{itemize}
  \end{block}
  \begin{exampleblock}{Na Parte 2}
    Dois algoritmos em que o invariante é menos óbvio --- e mais necessário:
    \textbf{busca binária} (famosa por implementações erradas) e o
    \textbf{particionamento} do quicksort.
  \end{exampleblock}
\end{frame}

\section*{Referências}
\begin{frame}[allowframebreaks]{Referências}
  \nocite{*}
  \bibliographystyle{plain}
  \bibliography{../referencias}
\end{frame}

\end{document}
